\hypertarget{appendix-73f5bcf7a462}{}\label{referencestart:appendix-73f5bcf7a462}
\section{Polar decomposition}\label{appendix-73f5bcf7a462-polar-decomposition}

The right polar form \(A=QH\) separates an orthogonal factor from a positive-semidefinite stretch \(H=(A^\mathsf T A)^{1/2}\). The orthogonal factor may include reflection; singular matrices can make its completion nonunique. Symmetric stretch can include shear in the chosen coordinates.

\textbf{Use it when:} separating local rigid-like motion from symmetric deformation.

\textbf{Developed in Phase 8:} Polar decomposition factors a map into an orthogonal part and a positive-semidefinite stretch. The SVD constructs both left and right forms, separates reflection from a proper rotation, and gives the nearest orthogonal or special-orthogonal factor when the singular directions allow a unique choice; zero singular values leave completion freedom. \href{../02-Mathematics-for-Visual-Computing/main.pdf\#nameddest=02-foundations-ii-ch024-polar-decomposition-separates-frame-dependent-stretch-from-orientation}{left and right stretch (Mathematics for Visual Computing, Mathematics for Visual Computing, Chapter 8)} · \href{../02-Mathematics-for-Visual-Computing/main.pdf\#nameddest=02-foundations-ii-ch024-two-sided-rotations-nearest-proper-rotation-and-nonunique-directions}{nearest proper rotation and ambiguity (Mathematics for Visual Computing, Mathematics for Visual Computing, Chapter 8)}

\subsection{Prerequisites}\label{appendix-73f5bcf7a462-prerequisites}

\hyperref[appendix-1fad5e7f36fa]{Singular value decomposition}, \hyperref[appendix-0b3be18411f9]{Eigenvalues and eigenspaces}, \hyperref[appendix-da478be01d02]{Determinants orientation and volume}.

\subsection{Used by}\label{appendix-73f5bcf7a462-used-by}

\hyperref[appendix-0f4bc2bf2b1f]{Matrix functions}, \hyperref[appendix-292c061cd47c]{Transform hierarchies TRS and shear}.

\textbf{Teaching path:} \href{../02-Mathematics-for-Visual-Computing/main.pdf\#nameddest=02-foundations-ii-ch024-polar-decomposition-separates-rotation-like-motion-from-stretch}{polar factors (Mathematics for Visual Computing, Mathematics for Visual Computing, Chapter 8)}.
